% =====================================================================
%  CARNET D'ENTRAINEMENT — FICHE 6 — THÈME 1
%  Niveau FACILE — Exercices
% =====================================================================
\documentclass[11pt,a4paper]{article}
\usepackage[utf8]{inputenc}
\usepackage[T1]{fontenc}
\usepackage{lmodern}
\usepackage[french]{babel}
\usepackage{amsmath,amssymb,mathtools}
\usepackage{icomma}
\usepackage[version=4]{mhchem}
\usepackage{siunitx}
\sisetup{output-decimal-marker={,}}
\usepackage{xcolor}
\usepackage{colortbl}
\usepackage{tikz}
\usepackage[most]{tcolorbox}
\usepackage{enumitem}
\usepackage{array}
\usepackage{booktabs}
\usepackage{fancyhdr}
\usepackage[hidelinks]{hyperref}
\usetikzlibrary{arrows.meta,positioning,calc,shapes.geometric}
\usepackage[a4paper,margin=2.1cm,headheight=26pt,headsep=8pt]{geometry}

\definecolor{primary}{RGB}{146,95,160}
\definecolor{primarydark}{RGB}{92,52,104}
\definecolor{excol}{RGB}{0,135,90}
\definecolor{atomO}{RGB}{220,55,45}
\definecolor{atomN}{RGB}{48,80,200}
\definecolor{atomH}{RGB}{235,235,235}
\definecolor{lonepair}{RGB}{120,94,200}
\definecolor{bondcol}{RGB}{60,60,70}

\tcbset{boxbase/.style={enhanced,breakable,boxrule=0.9pt,arc=2.5pt,
   left=3mm,right=3mm,top=2.2mm,bottom=2.2mm,
   fonttitle=\bfseries,coltitle=white,toptitle=0.6mm,bottomtitle=0.6mm}}
\newtcolorbox[auto counter]{exo}[1][]{boxbase,colback=primary!5,colframe=primary,
   title={\textsc{Exercice}~\thetcbcounter\ifx\relax#1\relax\relax\else\textmd{ —\itshape\ #1}\fi}}

\pagestyle{fancy}\fancyhf{}
\fancyhead[L]{\small\textcolor{primarydark}{\textbf{Fiche 6 · Thème 1} — Types de liaisons et octet}}
\fancyhead[R]{\small\textcolor{primarydark}{\textsc{Facile}}}
\fancyfoot[C]{\small\thepage}
\renewcommand{\headrulewidth}{0.8pt}
\renewcommand{\headrule}{\hbox to\headwidth{\color{primary}\leaders\hrule height \headrulewidth\hfill}}
\setlength{\parindent}{0pt}\setlength{\parskip}{3pt}

\begin{document}
\begin{center}
{\Large\bfseries\color{primarydark}Thème 1 — Niveau Facile}\\[2pt]
{\small\itshape Lire la fiche \texttt{Tuto\_Theme-1} avant de commencer. Aucune calculatrice n'est nécessaire.}
\end{center}
\medskip

\begin{exo}[électrons de valence]
En utilisant le numéro de colonne du tableau périodique, donner le nombre d'électrons de
\textbf{valence} de chacun de ces atomes :
\[ \ce{C}, \qquad \ce{O}, \qquad \ce{Cl}, \qquad \ce{H}, \qquad \ce{P}, \qquad \ce{Na}, \qquad \ce{S}, \qquad \ce{N}. \]
\end{exo}

\begin{exo}[duet ou octet ?]
Pour chacun des atomes suivants, indiquer de combien d'électrons il cherche à s'entourer lorsqu'il
se lie : \textbf{2} (duet) ou \textbf{8} (octet) ?
\[ \ce{H}, \qquad \ce{O}, \qquad \ce{C}, \qquad \ce{He}, \qquad \ce{F}, \qquad \ce{N}. \]
\end{exo}

\begin{exo}[nommer le type de liaison]
Indiquer, pour chaque situation, s'il s'agit d'une liaison \textbf{covalente}, \textbf{ionique}
ou \textbf{covalente de coordination (dative)} :
\begin{enumerate}[label=\alph*),leftmargin=2em,itemsep=1pt]
  \item la liaison entre les deux atomes de la molécule de dichlore \ce{Cl2} ;
  \item la cohésion dans le solide chlorure de sodium \ce{NaCl} ;
  \item la quatrième liaison \ce{N-H} formée lorsque \ce{NH3} capte un proton \ce{H+} ;
  \item la liaison dans le dihydrogène \ce{H2} (\ce{H-H}).
\end{enumerate}
\end{exo}

\begin{exo}[doublets liants et non liants]
On donne ci-dessous la structure de Lewis de la molécule d'eau \ce{H2O}. Autour de l'atome
d'oxygène, compter :
\begin{enumerate}[label=\alph*),leftmargin=2em,itemsep=1pt]
  \item le nombre de \textbf{doublets liants} ;
  \item le nombre de \textbf{doublets non liants} (paires libres).
\end{enumerate}

\begin{center}
\begin{tikzpicture}[scale=1.0]
  \coordinate (O) at (0,0);
  \coordinate (Ha) at (-1.3,-0.9);
  \coordinate (Hb) at ( 1.3,-0.9);
  \draw[bondcol,line width=1.4pt] (O)--($(Ha)+(0.25,0.17)$);
  \draw[bondcol,line width=1.4pt] (O)--($(Hb)+(-0.25,0.17)$);
  \node[font=\large\bfseries] at (O){O};
  \node[font=\large\bfseries] at (Ha){H};
  \node[font=\large\bfseries] at (Hb){H};
  \foreach \a in {60,120}{\fill[lonepair] ($(O)+(\a:0.55)$) circle (2.4pt);}
  \foreach \a in {90}{}
  \fill[lonepair] ($(O)+(75:0.8)$) circle (2.4pt);
  \fill[lonepair] ($(O)+(105:0.8)$) circle (2.4pt);
\end{tikzpicture}
\end{center}
\end{exo}

\begin{exo}[l'octet est-il respecté ?]
On donne la structure de Lewis de l'ammoniac \ce{NH3} : l'azote y forme trois liaisons \ce{N-H}
et porte un doublet non liant.
\begin{enumerate}[label=\alph*),leftmargin=2em,itemsep=1pt]
  \item Combien d'électrons entourent l'atome d'azote ? (compter les liaisons \emph{et} le doublet libre)
  \item L'azote respecte-t-il la règle de l'octet ?
\end{enumerate}
\end{exo}

\end{document}
